\chapter*{\IfLanguageName{dutch}{Samenvatting}{Abstract}}

Traditionele reanimatietraining voor verpleegkundigen steunt vaak op externe 2D-interfaces voor feedback, wat leidt tot het \textit{split-attention effect}. 
De hulpverlener moet hierbij de aandacht verdelen tussen de fysieke handeling op de reanimatiepop en een extern scherm, wat de cognitieve belasting verhoogt en de 
leercurve belemmert. Deze bachelorproef onderzoekt hoe Mixed Reality (MR) technologie dit proces kan optimaliseren door real-time, intuïtieve feedback direct in het 
gezichtsveld van de gebruiker te projecteren.

Het doel van dit onderzoek was de ontwikkeling van een Proof of Concept (PoC) op de Meta Quest 3, 
waarbij sensordata van een reanimatiepop in real-time wordt gevisualiseerd via een Unity-applicatie. 
Een significante technische barrière in dit proces was de gesloten aard van de beschikbare communicatieprotocollen.
Hoewel de fabrikant eigen ecosystemen aanbiedt, ontbreekt een publiek toegankelijke en gedocumenteerde SDK die directe 
integratie met de Meta Quest 3 en de Unity-omgeving mogelijk maakt. Daarnaast vormt de beperkte native ondersteuning voor Bluetooth Low Energy (BLE) 
binnen de standaard Unity-ontwikkelomgeving voor Android-gebaseerde headsets een extra hindernis.

Om dit probleem te overbruggen, is een diepgaand technisch onderzoek uitgevoerd naar het communicatieprotocol van de pop. 
Door middel van \textit{reverse engineering} technieken waaronder het decompileren van officiële software en het analyseren van 
Bluetooth Low Energy (BLE) pakketjes via \textit{HCI snoop logs} en Wireshark op een gerouteerd Android-toestel zijn de verborgen 
handshakes en data-characteristics ontcijferd. 

Het resultaat is een werkend prototype dat live sensordata (zoals compressiediepte en frequentie) vertaalt naar ruimtelijke 
visualisaties, waaronder een \textit{ghost avatar} en een dynamisch reagerend hartmodel. De voorlopige resultaten van deze PoC 
tonen aan dat MR in staat is om de kloof tussen haptische ervaring en visuele feedback te dichten. Dit onderzoek legt hiermee de 
technische fundamenten voor een nieuwe generatie medische simulatietools die verpleegkundigen effectiever kunnen voorbereiden op kritieke levensreddende situaties.

\textbf{Trefwoorden:} Mixed Reality, Meta Quest 3, Reanimatietraining, Reverse Engineering, Bluetooth Low Energy, Unity, Zorginnovatie.